%==============================================================
% CV Jorge Armando Cordero García
% VERSIÓN 1: Especialista en Soporte TI / Help Desk / Service Desk
% Optimizado para ATS | MiKTeX | VS Code
% Compilar con: pdflatex CV_Jorge_Cordero_SoporteTI.tex
%==============================================================

\documentclass[10pt, letterpaper]{article}

%------ PAQUETES ------
\usepackage[utf8]{inputenc}
\usepackage[T1]{fontenc}
\usepackage{babel}
\usepackage[left=1.8cm, right=1.8cm, top=1.5cm, bottom=1.5cm]{geometry}
\usepackage{lmodern}
% \usepackage{microtype}  % removed because not installed in current TinyTeX environment
\usepackage{parskip}
\usepackage{enumitem}
% \usepackage{titlesec}
\usepackage{xcolor}
\usepackage{hyperref}
\usepackage{amssymb}
% \usepackage{fontawesome5}
\usepackage{multicol}
\usepackage{array}
\usepackage{tabularx}
\usepackage{booktabs}
\usepackage{tikz}
\usepackage{pgffor}

%------ COLORES ------
\definecolor{primary}{HTML}{1A3A5C}
\definecolor{accent}{HTML}{2E86AB}
\definecolor{lightgray}{HTML}{F5F5F5}
\definecolor{darkgray}{HTML}{444444}
\definecolor{medgray}{HTML}{666666}
\definecolor{rulecol}{HTML}{2E86AB}

%------ HYPERLINKS ------
\hypersetup{
    colorlinks=true,
    urlcolor=accent,
    linkcolor=primary,
    pdfauthor={Jorge Armando Cordero García},
    pdftitle={CV - Especialista en Soporte TI},
    pdfsubject={Curriculum Vitae}
}

%------ FORMATO DE SECCIONES ------
% Se usa el formato de sección por defecto en este entorno.

%------ COMANDOS PERSONALIZADOS ------
\newcommand{\cvheader}[5]{%
  \begin{center}
    {\LARGE\bfseries\color{primary} #1}\\[4pt]
    {\large\color{accent} #2}\\[6pt]
    {\small\color{darkgray}
      Ubicación: #3 \quad
      Tel: #4 \quad
      Email: \href{mailto:#5}{#5}
    }\\[2pt]
    {\small\color{darkgray}
      LinkedIn: \href{https://linkedin.com/in/jacoga}{linkedin.com/in/jacoga}
    }
  \end{center}
  \vspace{2pt}
  \textcolor{rulecol}{\rule{\linewidth}{2pt}}
  \vspace{4pt}
}

\newcommand{\cventry}[5]{%
  \noindent
  \begin{tabularx}{\linewidth}{@{}X r@{}}
    {\bfseries\color{primary} #1} & {\small\color{medgray} #2} \\
    {\small\color{accent} #3} $\cdot$ {\small\color{darkgray}\itshape #4} & \\
  \end{tabularx}
  \vspace{2pt}
  #5
  \vspace{6pt}
}

\newcommand{\cvitem}{\item[\textcolor{accent}{\small$\blacktriangleright$}]}

\setlist[itemize]{
  leftmargin=14pt,
  itemsep=1pt,
  parsep=0pt,
  topsep=2pt,
  label=\textcolor{accent}{\small$\blacktriangleright$}
}

%------ SKILL BAR ------
\newcommand{\skillbar}[2]{%
  \noindent
  \begin{tikzpicture}[baseline]
    \fill[lightgray] (0,0) rectangle (4cm, 0.18cm);
    \fill[accent] (0,0) rectangle (#2*0.04cm, 0.18cm);
  \end{tikzpicture}
  \hspace{4pt}{\small #1}\\[3pt]
}

\pagestyle{empty}

%==============================================================
\begin{document}
%==============================================================

% ENCABEZADO
\cvheader
  {Jorge Armando Cordero García}
  {Ingeniero de Sistemas | Especialista en Soporte TI \& Service Desk}
  {Gamarra, Cesar – Colombia}
  {+57 317 819 5354}
  {iscuande\_1982@hotmail.com}

%------ PERFIL PROFESIONAL ------
\section{Perfil Profesional}

Ingeniero de Sistemas con más de \textbf{10 años de experiencia} en soporte técnico multinivel (Nivel 1, Nivel 2, Help Desk y Service Desk), gestión de infraestructura IT y administración de redes LAN/WAN en sectores de salud, educación y gobierno. Sólidas competencias en resolución de incidencias bajo metodología \textbf{ITIL}, administración de \textbf{Windows Server}, \textbf{Active Directory} y herramientas de soporte remoto (TeamViewer, AnyDesk). Certificado por \textbf{Google} en IT Support y Ciberseguridad. Disponibilidad inmediata para modalidad presencial, remota o híbrida.

%------ EXPERIENCIA PROFESIONAL ------
\section{Experiencia Profesional}

\cventry
  {Ingeniero de Sistemas – Consultor TI}
  {Ene 2024 – Presente}
  {Construcciones \& Soluciones del Río S.A.S. / Tictelco S.A.S. / Clientes Privados}
  {Aguachica, Cesar – Híbrida}
  {
    \begin{itemize}
      \item Soporte técnico Nivel 1 y Nivel 2 a usuarios finales en entornos Windows 10/11 y servidores locales.
      \item Administración y monitoreo de servidores, garantizando disponibilidad continua del servicio.
      \item Instalación, configuración y mantenimiento preventivo/correctivo de equipos de cómputo e impresoras.
      \item Diagnóstico y resolución de incidencias de conectividad en redes cableadas estructuradas (LAN).
      \item Gestión de racks de comunicaciones, switching y revisión de cableado estructurado.
      \item Soporte en plataformas WordPress, HTML/CSS y optimización de sistemas operativos.
      \item Capacitación técnica a usuarios sobre herramientas ofimáticas y buenas prácticas de seguridad.
    \end{itemize}
  }

\cventry
  {Ingeniero de Sistemas}
  {Feb 2023 – Jul 2024}
  {E.S.E. Hospital Olaya Herrera}
  {Gamarra, Cesar – Presencial}
  {
    \begin{itemize}
      \item Soporte TI multinivel a más de \textbf{200 usuarios} en entorno hospitalario de misión crítica.
      \item Administración de \textbf{Windows Server} y redes LAN; logro de \textbf{99,5\% de continuidad operativa}.
      \item Gestión y procesamiento de informes institucionales en plataformas SISPRO, MIPRES, RUAF y SIGIRES.
      \item Mantenimiento de servidores locales, equipos de cómputo, impresoras y activos tecnológicos.
      \item Configuración y administración de infraestructura de red LAN hospitalaria.
    \end{itemize}
  }

\cventry
  {Ingeniero de Sistemas}
  {Feb 2022 – Jul 2022}
  {Universidad Popular del César}
  {Aguachica, Cesar – Presencial}
  {
    \begin{itemize}
      \item Administración de infraestructura tecnológica institucional: servidores, redes y equipos de cómputo.
      \item Control de inventario TI y gestión de activos tecnológicos.
      \item Soporte técnico a usuarios y administración de sistemas académicos institucionales.
    \end{itemize}
  }

\cventry
  {Lider GELT – Gestión TIC}
  {Jun 2015 – Feb 2016}
  {Alcaldía Municipal de Gamarra}
  {Gamarra, Cesar – Presencial}
  {
    \begin{itemize}
      \item Liderazgo del área GELT (Gestión Electrónica) municipal; administración de servidores y red LAN.
      \item Gestión de recursos TIC e inventario del área tecnológica municipal.
    \end{itemize}
  }

\cventry
  {Ingeniero de Sistemas}
  {Mar 2016 – Mar 2017}
  {Linktech de Colombia}
  {Bogotá, D.C. – Presencial}
  {
    \begin{itemize}
      \item Mantenimiento preventivo y correctivo de equipos de cómputo; instalación de hardware/software.
      \item Administración de redes, cableado estructurado y certificación de infraestructura de red.
    \end{itemize}
  }

\cventry
  {Ingeniero de Sistemas}
  {Ene 2017 – Dic 2017}
  {Cordisco}
  {Bogotá, D.C. – Presencial}
  {
    \begin{itemize}
      \item Administración de red corporativa y soporte técnico interno a usuarios.
      \item Implementación de mejoras tecnológicas y gestión del Drive institucional.
    \end{itemize}
  }

%------ FORMACIÓN ACADÉMICA ------
\section{Formación Académica}

\noindent
\begin{tabularx}{\linewidth}{@{}X r@{}}
  {\bfseries\color{primary} Ingeniería de Sistemas} & {\small\color{medgray} 2015 – 2023} \\
  {\small\color{darkgray} Universidad Popular del César – Aguachica, Cesar} & \\
\end{tabularx}

\vspace{4pt}

\noindent
\begin{tabularx}{\linewidth}{@{}X r@{}}
  {\bfseries\color{primary} Especialización en TIC para Diseño de Estrategias Didácticas} & {\small\color{medgray} 2022 – 2023} \\
  {\small\color{darkgray} Universidad Popular del César – En proceso de grado} & \\
\end{tabularx}

%------ CERTIFICACIONES ------
\section{Certificaciones}

\begin{itemize}
  \item \textbf{Google IT Support Professional Certificate} – Google / Coursera (2025)
  \item \textbf{Google Data Analytics Professional Certificate} – Google / Coursera (2025)
  \item \textbf{Google Cybersecurity Professional Certificate} – Google / Coursera (En curso 2025)
  \item \textbf{IT Customer Support Basics} – Cisco Networking Academy (2025)
  \item \textbf{Introduction to IoT} – Cisco Networking Academy (2025)
  \item \textbf{Introduction to Cybersecurity} – Cisco Networking Academy (2025)
  \item \textbf{Fundamentos de IA para Negocios} – Google (2024)
  \item \textbf{Herramientas de IA Generativa} – Microsoft (2024)
  \item \textbf{HTML5/CSS3 Avanzado, JavaScript ES6+, React.js} – LinkedIn Learning (2024)
  \item \textbf{Inglés Técnico Niveles 1–3} – SENA Virtual (2021)
\end{itemize}

%------ HABILIDADES TÉCNICAS ------
\section{Habilidades Técnicas}

\begin{multicols}{2}
\noindent\textbf{\color{primary}Soporte \& Help Desk}\\
{\small Soporte Nivel 1/2, ITIL, Help Desk, Service Desk,\\
TeamViewer, AnyDesk, Diagnóstico de Incidencias}

\vspace{6pt}
\noindent\textbf{\color{primary}Sistemas Operativos}\\
{\small Windows 10/11, Windows Server 2016/2019,\\
Linux (fundamentos), Active Directory}

\vspace{6pt}
\noindent\textbf{\color{primary}Redes \& Conectividad}\\
{\small LAN/WAN, TCP/IP, Switching, Routing,\\
Cableado Estructurado, Rack Management}

\columnbreak

\noindent\textbf{\color{primary}Hardware \& Mantenimiento}\\
{\small Mantenimiento preventivo/correctivo PC,\\
Impresoras, Periféricos, Ensamblaje}

\vspace{6pt}
\noindent\textbf{\color{primary}Seguridad Informática}\\
{\small Hardening básico, Security Awareness,\\
Fundamentos de Ciberseguridad}

\vspace{6pt}
\noindent\textbf{\color{primary}Herramientas Ofimáticas}\\
{\small Microsoft Office Suite (Excel, Word, PowerPoint),\\
Google Workspace, WordPress}
\end{multicols}

%------ COMPETENCIAS BLANDAS ------
\section{Competencias Transversales}

{\small Resolución de problemas complejos $\cdot$ Comunicación asertiva con usuarios $\cdot$ Orientación al cliente $\cdot$ Pensamiento analítico $\cdot$ Trabajo bajo presión $\cdot$ Atención al detalle $\cdot$ Adaptabilidad $\cdot$ Trabajo en equipo $\cdot$ Gestión del tiempo}

%------ IDIOMAS ------
\section{Idiomas}

\noindent
\begin{tabularx}{\linewidth}{@{}l X@{}}
  \textbf{Español:} & Nativo \\
  \textbf{Inglés:} & Básico–Técnico (A2) – Lectura y escritura de documentación técnica \\
\end{tabularx}

%------ PIE ------
\vspace{8pt}
\begin{center}
\textcolor{medgray}{\small Disponibilidad inmediata $\cdot$ Modalidad presencial, remota o híbrida $\cdot$ Movilidad nacional}
\end{center}

\end{document}
